\chapter{Совместимость полного multiplicity-кадра с одиночной картой c0}
\label{ch:v8-baryon-c0-full-multiplicity-frame-single-map-compatibility-gate}

Пусть $V_0,V_1,V_2$ --- trace-ортонормированный комплексный базис трёх
допустимых старо-новых карт. Общий gauge-ковариантный dissipator задаётся
положительной матрицей Коссаковского
\begin{equation}
 \mathcal D_C(A)=\sum_{i,j=0}^{2}C_{ij}
 \left(V_iAV_j^*-\frac12\{V_j^*V_i,A\}\right),
 \qquad C\succeq0.
 \label{eq:v8-baryon-c0-multiplicity-kossakowski-generator}
\end{equation}
Полный изотропный Hom-кадр соответствует
\begin{equation}
 C_{\rm full}=I_3,
 \qquad \rank C_{\rm full}=3.
 \label{eq:v8-baryon-c0-full-frame-kossakowski-rank}
\end{equation}
Одна комплексная карта $V(z)=\sum_i z_iV_i$ соответствует чистой
ковариации
\begin{equation}
 C_z=zz^*,
 \qquad z\ne0,
 \qquad \rank C_z=1.
 \label{eq:v8-baryon-c0-single-map-kossakowski-rank}
\end{equation}

\section{Инвариант минимального Kraus-ранга}

Векторизация трёх стрелок даёт матрицу $W$ ранга три. Матрица Чоя
полностью положительной части имеет вид
\begin{equation}
 J_C=W C W^*,
 \qquad
 \rank J_{\rm full}=3,
 \qquad
 \rank J_z=1.
 \label{eq:v8-baryon-c0-choi-rank-obstruction}
\end{equation}
Следовательно, минимальные среды Стайнспринга имеют размеры
\begin{equation}
 d_{\rm env}^{\rm full}=3,
 \qquad d_{\rm env}^{\rm single}=1.
 \label{eq:v8-baryon-c0-stinespring-environment-ranks}
\end{equation}
Унитарная или изометрическая смена Kraus-базиса не меняет эти ранги, так
что операторное равенство невозможно:
\begin{equation}
 \mathcal D_{I_3}\ne\gamma\mathcal D_{zz^*}
 \quad\text{для всех }z\ne0,\ \gamma>0.
 \label{eq:v8-baryon-c0-full-frame-not-single-map}
\end{equation}

\section{Почему центральный тест не видит различия}

На базисе $P_{21}/\sqrt{21},P_{\mathrm n}/\sqrt2$ все три координатные
стрелки имеют одну и ту же форму $D_0$. Поэтому
\begin{equation}
 D_{\rm cen}(C)=\Tr(C)D_0,
 \qquad
 D_0=\begin{pmatrix}
 2/21&-\sqrt{42}/21\\
 -\sqrt{42}/21&1
 \end{pmatrix}.
 \label{eq:v8-baryon-c0-central-restriction-sees-only-trace}
\end{equation}
В частности, при $\lVert z\rVert=1$
\begin{equation}
 D_{\rm cen}(I_3)=3D_0
 =D_{\rm cen}(3zz^*).
 \label{eq:v8-baryon-c0-central-single-map-false-equivalence}
\end{equation}
Совпадение центральной щели и неподвижной алгебры поэтому не является
сертификатом равенства полных процессов.

Например, если одиночная карта выбрана вдоль $V_0$, то для проектора на
вторую старую линию получаем
\begin{equation}
 \mathcal D_{I_3}(P_{e_R})
 =2(P_{a_0}-P_{e_R}),
 \qquad
 \mathcal D_{V_0}(P_{e_R})=0.
 \label{eq:v8-baryon-c0-operator-witness-full-versus-single}
\end{equation}

Rank-one редукция возможна только после выбора чистого вектора среды
\begin{equation}
 |\psi_z\rangle=z_0|0\rangle+z_1|1\rangle+z_2|2\rangle,
 \qquad [z]\in\mathbb{CP}^2
 \quad(\mathbb{RP}^2\text{ в Real-срезе}).
 \label{eq:v8-baryon-c0-environment-pure-state-selector}
\end{equation}
Это ровно прежнее недостающее multiplicity-данное, перенесённое в среду.

Итоговый контракт равен
\begin{equation}
 \text{central compatibility: да},
 \qquad
 \text{operator/Kraus compatibility: нет}.
 \label{eq:v8-baryon-c0-full-single-compatibility-ledger}
\end{equation}

\begin{theorem}[Kraus-ранговый запрет одиночной интерпретации]
Полный изотропный multiplicity-frame имеет минимальный комплексный
Kraus-ранг три и не может быть преобразован в один комплексный коннектор.
Равенство центральных ограничений после перенормировки скорости не снимает
этот операторный запрет.
\end{theorem}

\begin{nogocriterion}[Пурификация не создаёт селектор]
Выбор чистого состояния трёхмерной среды действительно даёт rank-one
ковариацию, но требует точки $\mathbb{CP}^2$ или $\mathbb{RP}^2$ как нового
входа. Нельзя считать такой выбор следствием уникального следа $M_5$.
\end{nogocriterion}

\begin{protocol}\begin{enumerate}
\item комплексных connector channels полного кадра: \textbf{$3$};
\item Kraus/Choi rank полного кадра: \textbf{$3$};
\item Kraus/Choi rank одиночной карты: \textbf{$1$};
\item операторная эквивалентность: \textbf{нет};
\item центральная эквивалентность после rate-rescaling: \textbf{да};
\item минимальная среда полного кадра: \textbf{$3$};
\item чистое состояние среды выбрано: \textbf{нет};
\item одиночный $c_0$ выведен: \textbf{нет};
\item следующий гейт: \textbf{селектор чистого состояния multiplicity-среды}.
\end{enumerate}\end{protocol}

Машинный сертификат сохранён в
\path{s2t_v8_baryon_c0_full_multiplicity_frame_single_map_compatibility_gate_results.json}.